% ============================================================
%  PLANTILLA DE INFORME — Probabilidad y Estadística
%  Proyecto: TablasFrecuencia_ProbEstadistica (v1.0.1)
%  Compilar con: pdflatex plantilla_informe.tex  (2 veces, por el índice)
%  También funciona directo en Overleaf.
% ============================================================
\documentclass[11pt,letterpaper]{article}

% ---------- Idioma y codificación ----------
\usepackage[utf8]{inputenc}
\usepackage[T1]{fontenc}
% Nota: si tu instalación de TeX tiene el paquete de idioma español
% (texlive-lang-spanish / babel-spanish, ya incluido en Overleaf), podés
% activar la localización automática descomentando la línea de abajo:
% \usepackage[spanish,es-noshorthands]{babel}
% Sin babel, los acentos igual funcionan bien (gracias a inputenc+fontenc);
% solo se pierde la traducción automática de "Contents"/fechas, que ya
% dejamos resuelta manualmente más abajo.
\renewcommand{\contentsname}{Índice}
\makeatletter
\renewcommand{\today}{\number\day\ de \ifcase\month\or
  enero\or febrero\or marzo\or abril\or mayo\or junio\or
  julio\or agosto\or septiembre\or octubre\or noviembre\or diciembre\fi
  \ de \number\year}
\makeatother

% ---------- Layout ----------
\usepackage[margin=2.5cm]{geometry}
\usepackage{setspace}
\onehalfspacing

% ---------- Tablas y gráficas ----------
\usepackage{graphicx}
\usepackage{booktabs}
\usepackage{caption}
\usepackage{float}

% ---------- Color y estilo ----------
\usepackage{xcolor}
\definecolor{accent}{HTML}{4F46E5}
\usepackage{titlesec}
\titleformat{\section}
  {\normalfont\Large\bfseries\color{accent}}{\thesection}{1em}{}
\titleformat{\subsection}
  {\normalfont\large\bfseries\color{accent!80!black}}{\thesubsection}{1em}{}

% ---------- Encabezado / pie de página ----------
\usepackage{fancyhdr}
\pagestyle{fancy}
\fancyhf{}
\fancyhead[L]{\small Probabilidad y Estadística}
\fancyhead[R]{\small \leftmark}
\fancyfoot[C]{\thepage}
\renewcommand{\headrulewidth}{0.4pt}

% ---------- Matemáticas ----------
\usepackage{amsmath}
\usepackage{amssymb}

% ---------- Enlaces ----------
\usepackage[colorlinks=true,linkcolor=accent,urlcolor=accent,citecolor=accent]{hyperref}

% ============================================================
%  DATOS DEL INFORME — completar antes de compilar
% ============================================================
\newcommand{\tituloInforme}{Título del informe}
\newcommand{\nombreEstudiante}{Nombre y Apellido}
\newcommand{\numeroCuenta}{N.º de cuenta / carné}
\newcommand{\nombreMateria}{Probabilidad y Estadística}
\newcommand{\nombreProfesor}{Nombre del/la docente}
\newcommand{\semestre}{Cuarto semestre}
\newcommand{\fechaEntrega}{\today}

\begin{document}

% ---------------- PORTADA ----------------
\begin{titlepage}
    \centering
    \vspace*{1.5cm}
    {\Large \textbf{Universidad}}\\[0.3cm]
    {\large Facultad / Carrera}\\[2.5cm]

    {\color{accent}\rule{\linewidth}{1.2pt}}\\[0.4cm]
    {\Huge \bfseries \tituloInforme}\\[0.4cm]
    {\color{accent}\rule{\linewidth}{1.2pt}}\\[2cm]

    {\large \nombreMateria\ --- \semestre}\\[0.3cm]
    {\large Docente: \nombreProfesor}\\[2.5cm]

    {\large \textbf{Presentado por:}}\\
    {\large \nombreEstudiante}\\
    {\large \numeroCuenta}\\[2.5cm]

    \vfill
    {\large \fechaEntrega}
\end{titlepage}

\tableofcontents
\newpage

% ---------------- INTRODUCCIÓN ----------------
\section{Introducción}
Breve contexto del informe: qué problema estadístico se aborda, por qué es relevante
y qué se va a presentar en las siguientes secciones.

% ---------------- OBJETIVOS ----------------
\section{Objetivos}
\subsection{Objetivo general}
Redactar el objetivo general del trabajo.

\subsection{Objetivos específicos}
\begin{itemize}
    \item Primer objetivo específico.
    \item Segundo objetivo específico.
    \item Tercer objetivo específico.
\end{itemize}

% ---------------- MARCO TEÓRICO ----------------
\section{Marco teórico}
Definiciones relevantes: frecuencia absoluta ($f_i$), frecuencia relativa ($h_i = f_i/n$),
frecuencia relativa porcentual, frecuencia acumulada ($F_i$), etc. Citar bibliografía
con \verb|\cite{}| cuando corresponda \cite{ejemplo2024}.

% ---------------- DESARROLLO / METODOLOGÍA ----------------
\section{Desarrollo}

\subsection{Datos}
Descripción de los datos recolectados: fuente, tamaño de la muestra ($n$), tipo de
variable (cualitativa, discreta, continua).

\subsection{Tabla de frecuencias}
La tabla \ref{tab:frecuencias} fue generada con la app \texttt{tablas\_frecuencia.html}
del proyecto (carpeta \texttt{app/}), exportada como CSV y pegada aquí.

\begin{table}[H]
    \centering
    \caption{Tabla de distribución de frecuencias}
    \label{tab:frecuencias}
    \begin{tabular}{lcccc}
        \toprule
        \textbf{Categoría} & \textbf{$f_i$} & \textbf{$h_i$} & \textbf{$h_i\%$} & \textbf{$F_i$} \\
        \midrule
        A & 4 & 0.40 & 40.0\% & 4 \\
        B & 3 & 0.30 & 30.0\% & 7 \\
        C & 2 & 0.20 & 20.0\% & 9 \\
        D & 1 & 0.10 & 10.0\% & 10 \\
        \midrule
        \textbf{Total} & \textbf{10} & \textbf{1.00} & \textbf{100.0\%} & --- \\
        \bottomrule
    \end{tabular}
\end{table}

\subsection{Gráfica}
La figura \ref{fig:grafica} se exportó como PNG directamente desde la app.

\begin{figure}[H]
    \centering
    % Reemplazar por la imagen exportada, p. ej.:
    % \includegraphics[width=0.7\linewidth]{grafica_bar.png}
    \fbox{\parbox{0.7\linewidth}{\centering\vspace{2cm}Insertar \texttt{grafica\_bar.png} aquí\vspace{2cm}}}
    \caption{Gráfica de barras de la distribución de frecuencias}
    \label{fig:grafica}
\end{figure}

\subsection{Análisis}
Interpretación de los resultados: qué categoría es la moda, cómo se distribuyen los
datos, observaciones relevantes.

% ---------------- CONCLUSIONES ----------------
\section{Conclusiones}
Síntesis de los hallazgos principales y respuesta a los objetivos planteados.

% ---------------- REFERENCIAS ----------------
\begin{thebibliography}{9}
\bibitem{ejemplo2024}
Apellido, N. (2024). \textit{Título del libro o material}. Editorial.
\end{thebibliography}

\end{document}
